\chapter{实验结论}\label{part:conclusions}

\section{实验目标与课程要求完成情况}

本项目围绕“输出表达、证据位置与执行反馈对大模型任务质量与开销的影响”完成了从数据采集、匿名化清洗、统计推断到零依赖交互可视化系统的全流程建设。全部成果严格遵循《数据分析与可视化实践》题目二（选项 B：可视交互系统）的交付结构，核心完成情况如表~\ref{tab:completion} 所示。

\begin{table}[!htbp]
\centering
\caption{课程实验核心任务完成情况总览}\label{tab:completion}
\ReportTableStyle
\begin{tabularx}{\linewidth}{@{}>{\raggedright\arraybackslash}p{3.6cm} >{\raggedright\arraybackslash}X@{}}
\toprule
\textbf{任务模块} & \textbf{实测结果与交付证据} \\
\midrule
多任务采集与冻结 & 覆盖推理、文档问答与代码三类任务；仓库累计 \textbf{273,186 条}物理响应，主展示选用 \textbf{267,610 条}；405+ 原始制品留存 SHA-256 指纹 \\
清洗与 21 列匿名化 & 产出统一表 \texttt{all\_records.csv.gz}（248,586 行），题目哈希脱敏、465 条环境排除置 \texttt{null} 不填零 \\
\addlinespace[2pt]
低维与多维统计发现 & 完成 5,000 次按题聚类 Bootstrap，覆盖 2 条低维、2 条多维与 1 条异常反弹发现（完整指标见附录~\ref{app:stats}） \\
\addlinespace[2pt]
“知微”可视交互系统 & 纯标准库启动（\texttt{app/launch.py}）；具备五球入口、五屏叙事、31 策略矩阵联动、三页七图版、数据表弹窗与代码 Diff 抽屉 \\
自动化验证与回归 & \texttt{check\_coursework.py}（185 链接 0 错误）、\texttt{zhiwei.mjs}（29/29 通过）、\texttt{zhiwei-gate.mjs}（12/12 通过）、零严重无障碍违例 \\
\bottomrule
\end{tabularx}
\end{table}

\section{核心研究结论}

基于保存的离线实验记录与聚类配对区间检验，本项目得出以下五项实证结论（对应课程要求的“2 条低维 + 2 条多维 + 1 条异常”框架）及一项代码生成补充结论：

\paragraph{结论一（低维 1）：温和简洁表达在跨任务场景中普遍压缩输出并保留质量}在保留分步推理要求的前提下追加温和简洁约束（\texttt{BeConcise}），于 \texttt{Qwen2.5 / GSM8K} 上将平均输出从 236.7 Token 压缩至 185.8 Token（降幅 21.5\%）的同时维持 92.0\% 的高分；跨推理、文档问答与代码生成 17 个条件组全部实现输出缩短（降幅 11.4\%–57.5\%），15/17 组得分点估计持平或上升。

\paragraph{结论二（低维 2）：中文作答模板在 MATH500 上缩短输出并助三模型提分}在相同英文竞赛数学题上采用 \texttt{Respond in Chinese.} 模板后，四个受测模型全部实现 Token 节省（4.4\%–33.3\%），其中 \texttt{Qwen3}、\texttt{HY4} 与 \texttt{Qwen2.5} 分别提分 +11.5、+5.5 与 +2.5 个百分点，核心动因在于紧凑表达大幅降低了长推理触发 800 Token 上限截断的比例。

\paragraph{结论三（异常现象）：极端长度约束引发输出反弹与解题预算挤占}当词数上限从 10 压紧至 1 或字母上限从 100 压紧至 10 时，12 个模型 $\times$ 题库分组中各有 9 组出现实际返回 Token 反向增长。\texttt{DeepSeek / MATH500} 在 10 字母限制下输出暴涨 5 倍（102.0 $\to$ 520.6 Token）而得分暴跌 45.5 个百分点，说明过严格式约束迫使模型将预算消耗在字符计数与格式自纠上。

\paragraph{结论四（多维 1）：模型间的文档问答优势高度依赖证据所在位置}在 20 篇文档的 9 位置剖面实验中，\texttt{Qwen3} 相对 \texttt{Qwen2.5} 的优势在第 1 位高达 +16.50 个百分点（[9.50, 23.50]），但至第 20 位萎缩为 +2.50 个百分点（[-4.50, 9.00]），首尾增益落差达 14.00 pp（[5.50, 23.00]）。若脱离上下文位置仅看平均命中率，便会掩盖近四成题目（39.0\%）的动态翻转。

\paragraph{结论五（多维 2）：具体执行反例最有效地驱动多轮代码修复累积收益}在隔离隐藏强测的双层评测中，无反馈自检在四轮修改后改善微弱且出现回退，而直接提供公开诊断中的失败输入与期望/实际输出反例，使 \texttt{Qwen3-8B} 最终通过率从初稿 75.2\% 攀升至多反例组的 \textbf{85.6\%}（421/492，初错修复 51/122 且零回退），相较无反馈自检净高 +9.15 个百分点（[5.49, 13.41]）。

\section{解释边界与后续改进}

\textbf{相关性与因果边界}\quad 本系统呈现的全部趋势与置信区间均受限于当前固定的提示文本、受测模型版本及本地评测沙箱。我们不将接口返回的明文步骤等价于模型内部隐式推理过程，也不将各任务独立计算的原生指标强行混合为跨任务“总准确率”。此外，由于文档位置实验的两个模型采集日期不同，且相同接口别名不保证服务端权重永久一致，结论仅对当前归档快照有效。

\textbf{后续演进方向}\quad 系统架构上可进一步补充两点能力：一是打通浏览器端对用户自定义 JSONL 对话日志的本地拖拽解析与实时 21 列 Schema 预览；二是将代码修复补充实验中的失败反例诊断模式接入真实 IDE 插件，进而构建具备预算熔断提醒的轻量化提示词辅助面板。


\textbf{实践工程启示}\quad 从本次全流程系统研发来看，大模型行为评估必须坚持“三不原则”：不合并异构任务的原始评分标尺，不用零值掩盖缺失与环境排除状态，不在缺乏多位置或多轮迭代切片时轻易下全局排名结论。交互系统的价值正是在于把这些边界条件显式交还给分析者自行核验。
